\clearpage\fancyhead[L]{\small COMET.}\SourceLeaf{350}
\begin{center}\textbf{COMET.}\par\vspace{4pt}\hrule\vspace{4pt}\fontsize{10}{12}\selectfont Monatliche Correspondenz. Band XXV. Seite 94..97. 1812. Januar.\par\vspace{4pt}\hrule\vspace{4pt}\end{center}

\par\indent [Der Inhalt dieser Briefe ist bis auf den letzten Absatz auch in dem Artikel vom 25. Januar 1812 der Gött. gel. Anz. wiedergegeben, nur steht hier Perihel Zeit .\;.\;.\;.\;. 1811 Nov. 11,392108 Merid. von Göttingen.]

\par\indent Die Bedeckung von $\lambda$ Aquarii hat Hr. \textsc{Wachter} neuerlich auch berechnet und folgende verbesserte Conjunctionszeiten gefunden:

\begin{center}
\fontsize{10}{12}\selectfont\setlength{\tabcolsep}{6pt}\renewcommand{\arraystretch}{1.15}\begin{tabular}{l|ccc}
Göttingen & $10^{\text{u}}$ & $22^{\text{m}}$ & $29^{\text{s}}69$ \\
Mannheim & 10 & 16 & 38.01 \\
Cronach & 10 & 28 & 3.50 \\
\end{tabular}
\end{center}

\noindent Die beiden ersten stimmen gut mit Ihnen, aber bei der letzten ist ein Unterschied von 10 Secunden. Da er die Rechnung doppelt nach zwei verschiedenen Methoden geführt hat, so halten Sie es vielleicht der Mühe nicht unwerth, Ihre letzte Rechnung noch einmal nachzusehen.
